\documentclass[10pt,a4paper]{amsart}
\usepackage[margin=19mm]{geometry}
\usepackage{amsmath,amssymb,mathtools}
\usepackage[scr=rsfs,cal=euler]{mathalfa}
\usepackage{xfrac}
\usepackage[unicode,hidelinks,bookmarks=false]{hyperref}
\newcommand{\ie}{i.e.\ }
\newcommand{\cf}{cf.\ }
\newcommand{\resp}{respectively\ }
\newcommand{\Subsection}{\subsection}
\providecommand{\nobreakdash}{\nobreak\hbox{-}}
\setlength{\emergencystretch}{2em}
\multlinegap0pt
\usepackage{bm}
\usepackage{bbm}
\usepackage{dsfont}
\usepackage{xspace}
\usepackage{cancel}
\usepackage{mathtools}
\usepackage{amsthm}
\makeatletter
\@ifundefined{thm}{\newtheorem{thm}{Theorem}}{}
\makeatother
\newcommand{\ds}{\displaystyle}
\title[Asymptotic behavior for a finitely degenerate semilinear pse]{Asymptotic behavior for a finitely degenerate semilinear pseudo-parabolic equation}
\author{Xiang-kun Shao, Xue-song Li, Nan-jing Huang, Donal O'Regan}
\date{Source: arXiv:2411.12253v2 (CC BY 4.0)}
\begin{document}
\maketitle

\noindent\emph{Source and licence.} Asymptotic behavior for a finitely degenerate semilinear pseudo-parabolic equation. Version arXiv:2411.12253v2 (CC BY 4.0), distributed under the Creative Commons Attribution 4.0 International licence, as recorded in the arXiv licence metadata for this exact version and shown on the abstract page https://arxiv.org/abs/2411.12253v2. Full rights notice: licenses/arxiv-2411.12253-CC-BY-4.0.txt.\par\smallskip

\noindent\textit{Editorial scope.} Counted unit: the Thm whose source label is \texttt{xx-33} in the article (source0.tex lines 606--611, source label \texttt{xx-33}), delivered together with the article's own complete original proof of it (source0.tex lines 614--777, one \texttt{proof} environment of 164 lines). No mathematics is written, repaired or re-derived: statement and proof are the article's own text line for line. Required context retained uncounted: 1.1 (lines 61--67); Hfanshu (lines 121--123); tezhengzhi (lines 125--127); x1111-33 (lines 130--132); J (lines 136--138); JI (lines 144--146); qianru (lines 160--162); weakdengshi (lines 172--174); nengliangdengshi (lines 180--182); baopo (lines 296--349); 68 (lines 328--330). Every definition, display and label of those lines is retained unchanged. Omitted from this package and not redistributed: 1--60, 68--120, 124--124, 128--129, 133--135, 139--143, 147--159, 163--171, 175--179, 183--295, 331--605, 612--613, 778--1027. Nothing inside a retained excerpt is omitted or altered: every retained body is delivered exactly as the article's lines are written, so no omission receipt of any kind accompanies this package. Citations: the retained excerpts carry no citation command at all, so no bibliography entry of the article is transcribed and no reference is redistributed. Reference adaptation: none was needed and none was made; every reference inside the delivered unit resolves inside it, so no unresolved reference or citation is produced. Printed slips kept literal: no printed word, symbol, hypothesis or number of the counted statement or of its proof was altered. Layout and class: the article's own class is \texttt{article} with packages including latexsym, amsfonts, amsthm, amsmath, amssymb, cases, cite, enumerate, graphicx, color, bm; the wrapper is the lane's pinned \texttt{amsart} editorial wrapper at 10pt on A4, which restates the theorem-like environments the retained text needs and adds the article's own macro definitions that the retained text needs (\texttt{\textbackslash ds}), with extra wrapper packages (bm, bbm, dsfont, xspace, cancel, mathtools). The delivered numbering is the unit's own. Assets: no retained span contains a figure environment, a graphics-inclusion command, a PDF-page inclusion, a table or a TikZ picture; no image file is copied into this package, the package declares no asset dependency and no image file is redistributed with it. Licence: version arXiv:2411.12253v2 of this article is distributed under the Creative Commons Attribution 4.0 International licence, as recorded in the arXiv licence metadata for this exact version and shown on the abstract page https://arxiv.org/abs/2411.12253v2. A full rights notice is delivered at licenses/arxiv-2411.12253-CC-BY-4.0.txt. Characters: every retained excerpt is pure ASCII, so the delivered engine, XeLaTeX with its default Latin Modern text font, reports no missing character.
\begin{align}\label{1.1}
 \left\{\begin{array}{ll}
     \ds \phi_{t}-\Delta_X\phi-\Delta_X\phi_t=|\phi|^{p-1}\phi,\ \ \ &x\in\Omega,\ t>0,\\
     \ds \phi(x,t)=0,&x\in\partial\Omega,\ t>0,\\
     \ds \phi(x,0)=\phi_0(x),\ \ &x\in\Omega,
    \end{array}\right.
\end{align}

\begin{equation}\label{Hfanshu}
\|\varphi\|_{H_{X,0}^1(\Omega)}^2=\|X\varphi\|_2^2+\|\varphi\|_2^2,\quad\forall\varphi\in H_{X,0}^1(\Omega).
\end{equation}

\begin{equation}\label{tezhengzhi}
\lambda_1\|\varphi\|_2^2\leq\|X\varphi\|_2^2,\ \ \ \forall\varphi\in H_{X,0}^1(\Omega).
\end{equation}

\begin{equation}\label{x1111-33}
\|\varphi\|_{p+1}\leq C\|\varphi\|_{H_{X,0}^1(\Omega)},\ \ \ \forall\varphi\in H_{X,0}^1(\Omega),
\end{equation}

\begin{align}
J(\varphi):=\frac{1}{2}\|X\varphi\|_2^2-\frac{1}{p+1}\|\varphi\|_{p+1}^{p+1}\label{J}
\end{align}

\begin{equation}\label{JI}
J(\varphi)=\frac{1}{p+1}I(\varphi)+\frac{p-1}{2(p+1)}\|X\varphi\|_2^2.
\end{equation}

\begin{equation}\label{qianru}
C_*:=\sup_{\varphi\in H_{X,0}^1(\Omega)\setminus\{0\}}\frac{\|\varphi\|_{p+1}}{\|X\varphi\|_2}.
\end{equation}

\begin{equation}\label{weakdengshi}
(\phi_t,\psi)+(X\phi,X\psi)+(X\phi_t,X\psi)=(|\phi|^{p-1}\phi,\psi),\quad  \forall\psi\in H_{X,0}^1(\Omega),~t\in(0,T).
\end{equation}

\begin{equation}\label{nengliangdengshi}
J(\phi)+\int_0^t\|\phi_\tau\|_{H_{X,0}^1(\Omega)}^2d\tau=J(\phi_0),\quad \forall t\in[0,T)
\end{equation}

\begin{thm}\label{baopo}
Suppose that $(A_1)$-$(A_3)$ hold. Let $\phi=\phi(t)$, $t\in[0,T)$ be a weak solution to \eqref{1.1}. If $J(\phi_0)<d$ and $I(\phi_0)<0$, then $\phi$ blows up at some finite time $T$, namely,
\begin{equation*}
\lim_{t\rightarrow T^-}\|\phi\|_{H_{X,0}^1(\Omega)}^2=\infty.
\end{equation*}
Furthermore,
\begin{description}
  \item{$(i)$} the upper bounds of blow-up time and rate are given by
  \begin{equation}\label{ut}
  T\leq\frac{\|u_0\|_{H_{X,0}^1(\Omega)}^2}{C_1(C_1-2)(d-J(\phi_0))},\hbox{~~~if~~}0\leq J(\phi_0)<d,
  \end{equation}
  and
  \begin{equation}\label{ur}
  \|\phi\|_{H_{X,0}^1(\Omega)}^2\leq\left\{
               \begin{array}{ll}
               \ds  \left[\frac{-C_1J(\phi_0)(C_1-2)}{\|\phi_0\|_{H_{X,0}^1(\Omega)}^{C_1}}\right]^{\frac{2}{2-C_1}}(T-t)^{-\frac{2}{C_1-2}},&\hbox{~if~~}J(\phi_0)<0; \\
               \vspace{-0.05in}\\
               \ds  \left[\frac{C_1(C_1-2)(d-J(\phi_0))}{\|\phi_0\|_{H_{X,0}^1(\Omega)}^{C_1}}\right]^{\frac{2}{2-C_1}}(T-t)^{-\frac{2}{C_1-2}},&\hbox{~if~~}0\leq J(\phi_0)<d,
               \end{array}
             \right.
  \end{equation}
  where
  \begin{equation}\label{C1}
  C_1=\left\{
               \begin{array}{ll}
               \ds  p+1,&\hbox{ if~~}J(\phi_0)<0; \\
               \vspace{-0.05in}\\
               \ds  \frac{\left(\epsilon_0^{p+1}-1\right)(p-1)}{\epsilon_0^{p+1}}+2,&\hbox{ if~~}0\leq J(\phi_0)<d.
               \end{array}
             \right.
  \end{equation}
  Here
  \begin{equation}\label{68}
  \epsilon_0=\left[\frac{p+1}{2}-C_*^{\frac{2(p+1)}{p-1}}J(\phi_0)(p+1)\right]^{\frac{1}{p-1}}>1,
  \end{equation}
  and $C_*>0$ is given in \eqref{qianru}.
  \item{$(ii)$} the exponential growth estimate of blow-up solutions is given by
  \begin{equation*}
  \|\phi\|_{H_{X,0}^1(\Omega)}^2\geq\left\{
               \begin{array}{ll}
               \ds  \|\phi_0\|_{H_{X,0}^1(\Omega)}^2e^{C_2t},&\hbox{ if }J(\phi_0)<0; \\
               \vspace{-0.05in}\\
               \ds  \|\phi_0\|_{H_{X,0}^1(\Omega)}^2e^{C_3t},&\hbox{ if }0\leq J(\phi_0)<d,
               \end{array}
             \right.
  \end{equation*}
  where
  \begin{equation}\label{C23}
  C_2=\frac{\lambda_1(p-1)}{\lambda_1+1},\quad
  C_3=\frac{\lambda_1(p-1)\left(\epsilon_0^2-1\right)}{\epsilon_0^2(\lambda_1+1)}.
  \end{equation}
  Here $\lambda_1$ and $\epsilon_0$ are given in \eqref{tezhengzhi} and \eqref{68}, respectively.
\end{description}
\end{thm}

\begin{thm}\label{xx-33}
Suppose that $(A_1)$-$(A_3)$ hold. Let $\phi=\phi(t)$ be a global solution to \eqref{1.1}. Then there exists $\phi^*\in \Psi$ and an increasing sequence $\{t_k\}_{k=1}^\infty$ with $t_k\rightarrow\infty$ as $k\rightarrow\infty$ such that
\begin{equation*}
\lim_{k\rightarrow\infty}\|\phi(t_k)-\phi^*\|_{H_{X,0}^1(\Omega)}=0.
\end{equation*}
\end{thm}

\begin{proof}
Let $\phi=\phi(t)$ be the global solution to problem \eqref{1.1}.
First, we prove
\begin{equation}\label{x1-33}
0\leq J(\phi)\leq J(\phi_0),\quad \forall t\in[0,\infty).
\end{equation}
From \eqref{nengliangdengshi}, it follows that $J(\phi(t))$ is nonincreasing with respect to $t$ and $J(\phi)\leq J(\phi_0)$ holds for all $t\in[0,\infty)$. We claim that $J(\phi)\geq 0$ holds for all $t\in[0,\infty)$. If not, there would exist $t_0\in[0,\infty)$ such that $J(\phi(t_0))<0$. Combining this with \eqref{JI}, we have $I(\phi(t_0))<0$. According to Theorem \ref{baopo}, $\phi$ would blow up in finite time, which contradicts the fact that $\phi$ is a global solution to problem \eqref{1.1}. Therefore, \eqref{x1-33} holds.

From \eqref{x1-33} and the fact that $J(\phi(t))$ is non-increasing with respect to $t$ on $[0,\infty)$, it follows that there exists a positive constant $C_0\in[0,J(\phi_0)]$ such that
\begin{equation}\label{x2-33}
\lim\limits_{t\rightarrow\infty}J(\phi)=C_0.
\end{equation}
Letting $t\rightarrow\infty$ in \eqref{nengliangdengshi}, we obtain
\begin{equation}\label{x3-33}
\int_0^\infty\|\phi'(\tau)\|_{H_{X,0}^1(\Omega)}^2d\tau=J(\phi_0)-C_0\leq J(\phi_0),
\end{equation}
which implies the existence of an increasing sequence $\{t_k\}_{k=1}^\infty$ with $t_k\rightarrow\infty)$ as $k\rightarrow\infty$ such that
\begin{equation}\label{x4-33}
\lim_{k\rightarrow\infty}\|\phi'(t_k)\|_{H_{X,0}^1(\Omega)}=0.
\end{equation}
According to \eqref{J} and \eqref{weakdengshi}, for any $\varphi\in H_{X,0}^1(\Omega)$, we have
\begin{equation}\label{x5-33}
\begin{split}
\langle J'(\phi(t)),\varphi\rangle&=\frac{d}{d\tau}J(\phi(t)+\tau\varphi)\bigg|_{\tau=0}\\
&=\left[\left(X(\phi(t)+\tau\varphi),X\varphi\right)-\left(|\phi(t)+\tau\varphi|^{p-1}(\phi(t)+\tau\varphi),\varphi\right)\right]\big|_{\tau=0}\\
&=(X\phi(t),X\varphi)-\left(|\phi(t)|^{p-1}\phi(t),\varphi\right)\\
&=-(\phi'(t),\varphi)-(X\phi'(t),X\varphi).
\end{split}
\end{equation}
Using \eqref{x5-33}, \eqref{Hfanshu} and \eqref{tezhengzhi}, we deduce
\begin{equation*}
\begin{split}
\|J'(\phi(t_k))\|_{H_X^{-1}(\Omega)}&=\sup_{\|\varphi\|_{H_{X,0}^1(\Omega)}\leq1}\big|\langle J'(\phi(t_k)),\varphi\rangle\big|\\
&\leq\sup_{\|\varphi\|_{H_{X,0}^1(\Omega)}\leq1}\left(\|\phi'(t_k)\|_2\|\varphi\|_2+\|X\phi'(t_k)\|_2\|X\varphi\|_2\right)\\
&\leq\sup_{\|\varphi\|_{H_{X,0}^1(\Omega)}\leq1}\|X\varphi\|_2\left(\lambda_1^{-\frac{1}{2}}\|\phi'(t_k)\|_2+\|X\phi'(t_k)\|_2\right)\\
&\leq\left(\lambda_1^{-\frac{1}{2}}+1\right)\|\phi'(t_k)\|_{H_{X,0}^1(\Omega)},
\end{split}
\end{equation*}
which, combined with \eqref{x4-33}, yields
\begin{equation}\label{x6-33}
\lim_{k\rightarrow\infty}\|J'(\phi(t_k))\|_{H_X^{-1}(\Omega)}=0.
\end{equation}
So there exists a constant $c>0$ such that
\begin{equation*}
\begin{split}
|I(\phi(t_k))|&=|\langle J'(\phi(t_k)),\phi(t_k)\rangle|\\
&\leq\|J'(\phi(t_k))\|_{H_X^{-1}(\Omega)}\|\phi(t_k)\|_{H_{X,0}^1(\Omega)}\\
&\leq c\|\phi(t_k\|_{H_{X,0}^1(\Omega)}.
\end{split}
\end{equation*}
It follows from \eqref{x4-33}, \eqref{Hfanshu} and \eqref{tezhengzhi} that
\begin{equation*}
\begin{split}
J(\phi_0)+\frac{c}{p+1}\|\phi(t_k)\|_{H_{X,0}^1(\Omega)}&\geq J(\phi(t_k))-\frac{1}{p+1}I(\phi(t_k))\\
&=\frac{p-1}{2(p+1)}\|X\phi(t_k)\|_2^2\\
&\geq\frac{(p-1)\lambda_1}{2(p+1)(1+\lambda_1)}\|\phi(t_k)\|_{H_{X,0}^1(\Omega)}^2.
\end{split}
\end{equation*}
Thus, there exists a constant $K_1>0$ such that
\begin{equation}\label{x8-33}
\|\phi(t_k)\|_{H_{X,0}^1(\Omega)}\leq K_1, \quad k=1,2,\cdots.
\end{equation}
Since the embedding $H_{X,0}^1(\Omega)\hookrightarrow L^{p+1}(\Omega)$ is compact, there exists an increasing subsequence of $\{t_k\}_{k=1}^\infty$ (still denoted as $\{t_k\}_{k=1}^\infty$) and $u_*\in H_{X,0}^1(\Omega)$ such that, as $k\rightarrow\infty$ (denote $\phi_k:=\phi(t_k)$):
\begin{align}
&\phi_k\rightharpoonup \phi_* \hbox{ weakly in } H_{X,0}^1(\Omega),\label{x9-33}\\
&\phi_k\rightarrow \phi_* \hbox{ strongly in } L^{p+1}(\Omega).\label{x10-33}
\end{align}
By \eqref{J}, we have
\begin{equation*}
\begin{split}
\quad\langle J'(\phi_k),\phi_k-\phi_*\rangle&=\frac{d}{d\tau}J(\phi_k+\tau (\phi_k-\phi_*))\bigg|_{\tau=0}\\
&=\big\{\left(X[\phi_k+\tau(\phi_k-\phi_*)],X(\phi_k-\phi_*)\right)\\
&\quad-\left(|\phi_k+\tau(\phi_k-\phi_*)|^{p-1}[\phi_k+\tau(\phi_k-\phi_*)],\phi_k-\phi_*\right)\big\}\big|_{\tau=0}\\
&=\left(X\phi_k,X(\phi_k-\phi_*)\right)-\left(|\phi_k|^{p-1}\phi_k,\phi_k-\phi_*\right).
\end{split}
\end{equation*}
Similarly,
\begin{equation}\label{x12-33}
\langle J'(\phi_*),\phi_k-u_*\rangle=(X\phi_*,X(\phi_k-\phi_*))-\left(|\phi_*|^{p-1}\phi_*,\phi_k-\phi_*\right).
\end{equation}
Thus, we obtain
\begin{equation}\label{x16-33}
\quad \langle J'(\phi_k)-J'(\phi_*),\phi_k-\phi_*\rangle=\|X(\phi_k-\phi_*)\|_2^2-\Theta,
\end{equation}
where $$\Theta:=\left(|\phi_k|^{p-1}\phi_k-|\phi_*|^{p-1}\phi_*,\phi_k-\phi_*\right).$$
Then by H\"{o}lder inequality, \eqref{x1111-33} and \eqref{x8-33}, one has
\begin{equation*}
\begin{split}
|\Theta|&\leq\|\phi_*^p-\phi_k^p\|_{\frac{p+1}{p}}\|\phi_k-\phi_*\|_{p+1}\\
&\leq\left(\|\phi_*\|_{p+1}^p+\|\phi_k\|_{p+1}^p\right)\|\phi_k-\phi_*\|_{p+1}\\
&\leq\left(\|\phi_*\|_{p+1}^p+C^p\|\phi_k\|_{H_{X,0}^1(\Omega)}^p\right)\|\phi_k-\phi_*\|_{p+1}\\
&\leq\left(\|\phi_*\|_{p+1}^p+C^pK_1^p\right)\|\phi_k-\phi_*\|_{p+1},
\end{split}
\end{equation*}
which, combined with \eqref{x10-33}, implies
\begin{equation}\label{x11-33}
\lim_{k\rightarrow\infty}|\Theta|=0.
\end{equation}
From \eqref{x9-33}, \eqref{x10-33} and \eqref{x12-33}, we obtain
\begin{equation}\label{x13-33}
\lim_{k\rightarrow\infty}\langle J'(\phi_*),\phi_k-\phi_*\rangle=0.
\end{equation}
By \eqref{x8-33}, we have
\begin{equation*}
\begin{split}
\big|\langle J'(\phi_k),\phi_k-u_*\rangle\big|&\leq\|J'(\phi_k)\|_{H_X^{-1}(\Omega)}\left(\|\phi_k\|_{H_{X,0}^1(\Omega)}+\|\phi_*\|_{H_{X,0}^1(\Omega)}\right)\\
&\leq\|J'(\phi_k)\|_{H_X^{-1}(\Omega)}\left(K_1+\|\phi_*\|_{H_{X,0}^1(\Omega)}\right),
\end{split}
\end{equation*}
and combining this with \eqref{x6-33}, we arrive at
\begin{equation}\label{x15-33}
\lim_{k\rightarrow\infty}\big|\langle J'(\phi_k),\phi_k-\phi_*\rangle\big|=0.
\end{equation}
It follows from \eqref{x16-33}, \eqref{x11-33}, \eqref{x13-33} and \eqref{x15-33} that
\begin{equation*}
\lim_{k\rightarrow\infty}\|X(\phi_k-\phi_*)\|_2^2=\lim_{k\rightarrow\infty}\left[\langle J'(\phi_k)-J'(\phi_*),\phi_k-\phi_*\rangle+\Theta\right]=0.
\end{equation*}
By \eqref{Hfanshu} and \eqref{tezhengzhi}, we conclude
\begin{equation*}
\lim_{k\rightarrow\infty}\|\phi_k-\phi_*\|_{H_{X,0}^1(\Omega)}^2=0.
\end{equation*}

Finally, we prove $\phi_*\in \Psi$. From \eqref{x10-33}, there exists a subsequence of $\{\phi_k\}_{k=1}^\infty$ (still denoted as $\{\phi_k\}_{k=1}^\infty$) such that, as $k\rightarrow\infty$,
\begin{equation}\label{x21-33}
\phi_k(x)\rightarrow \phi_*(x) \hbox{ a.e. in } \Omega,
\end{equation}
and there exists $w\in L^{p+1}(\Omega)$ satisfying, for all $k$,
\begin{equation}\label{x22-33}
|\phi_k(x)|\leq w(x) \hbox{ a.e. in } \Omega.
\end{equation}
By \eqref{x21-33}, we conclude that, as $k\rightarrow\infty$,
\begin{equation}\label{x24-33}
|\phi_k(x)|^{p-1}\phi_k(x)\nu(x)\rightarrow |\phi_*(x)|^{p-1}\phi_*(x)\nu(x) \hbox{ a.e. in } \Omega
\end{equation}
for $\nu\in H_{X,0}^1(\Omega)$, and it follows from \eqref{x22-33} that, for all $k$,
\begin{equation}\label{x25-33}
\left||\phi_k(x)|^{p-1}\phi_k(x)\nu(x)\right|\leq \left||w(x)|^{p-1}w(x)\nu(x)\right| \hbox{ a.e. in } \Omega.
\end{equation}
Applying H\"{o}lder inequality and \eqref{x1111-33}, we obtain
\begin{equation*}
\int_\Omega\left||w(x)|^{p-1}w(x)\nu(x)\right|dx\leq\left\||w(x)|^p\right\|_{\frac{p+1}{p}}\|\nu(x)\|_{p+1}\leq C\|w(x)\|_{p+1}^p\|\nu(x)\|_{H_{X,0}^1(\Omega)}.
\end{equation*}
Combining $w\in L^{p+1}(\Omega)$ and $\nu\in H_{X,0}^1(\Omega)$, one has
\begin{equation*}
\int_\Omega\left||w(x)|^{p-1}w(x)\nu(x)\right|dx<\infty.
\end{equation*}
Hence,
\begin{equation}\label{vvv}
|w(x)|^{p-1}w(x)\nu(x)\in L^1(\Omega).
\end{equation}
By the Lebesgue dominated convergence theorem, together with \eqref{x24-33}, \eqref{x25-33} and \eqref{vvv}, it follows that, for any $\nu\in H_{X,0}^1(\Omega)$,
\begin{equation}\label{x23-33}
\left(|\phi_k|^{p-1}\phi_k,\nu\right)\rightarrow\left(|\phi_*|^{p-1}\phi_*,\nu\right) \mbox{  as  } k\rightarrow\infty.
\end{equation}
Since
\begin{equation*}
\langle J'(\phi_k),\nu\rangle=(X\phi_k,X\nu)-\left(|\phi_k|^{p-1}\phi_k,\nu\right),
\end{equation*}
letting $k\rightarrow \infty$ and using \eqref{x6-33}, \eqref{x9-33} and \eqref{x23-33}, we obtain
\begin{equation*}
0=(X\phi_*,X\nu)-\left(|\phi_*|^{p-1}\phi_*,\nu\right).
\end{equation*}
Therefore, $\phi_*\in \Psi$.
\end{proof}

\end{document}
